\section{方法}\label{sec:method}
当前实现的 CAME-Net 由多模态射影嵌入（MPE）、几何 Clifford 注意力（GCA）、几何层归一化（GLN）以及软等变正则化组成。点云、图像和文本首先分别映射到 $\mathrm{Cl}(3,0,1)$ 的统一多重向量空间，随后沿 token 维拼接并送入共享的 GCA 堆栈，最后通过全局平均池化和标准分类头输出任务结果。

本文并不声称端到端网络在解析意义下严格满足 $\mathrm{SE}(3)$ 等变性；更准确地说，GCA 的标量打分机制构成严格几何核心，而 GLN 与按阶参数化带来端到端近似误差。因此，本文方法的恰当表述是：具有严格几何打分核心、并在端到端意义上近似 $\mathrm{SE}(3)$ 等变的多模态网络。

\subsection{多模态射影嵌入}\label{sec:mpe}
给定点云坐标 $\{\mathbf{p}_i\}_{i=1}^N$，实现中先通过 $K$ 近邻协方差分解估计局部法向 $\mathbf{n}_i$，再构造刚体运动不变量描述子并用 MLP 生成标量通道 $s_i\in\langle\mathrm{Cl}\rangle_0$。点云 token 写为
\begin{equation}\label{eq:mpe_pc_impl}
\mathcal{X}_i^{\mathrm{pc}} = s_i + \Pi_i + B_i + P(\mathbf{p}_i),
\end{equation}
其中 $\Pi_i = n_{ix}e_1 + n_{iy}e_2 - n_{iz}e_3 - (\mathbf{n}_i^\top\mathbf{p}_i)e_0$ 为局部切平面，$B_i = n_{ix}e_{23} + n_{iy}e_{31} + n_{iz}e_{12}$ 为欧几里得二向量，且显式排除了理想二向量 $\{e_{01},e_{02},e_{03}\}$；$P(\mathbf{p}_i)=e_{123}+x_ie_{032}+y_ie_{013}+z_ie_{021}$ 为点的标准 OPNS PGA 表示。

图像与文本分支只激活标量与伪标量通道，即
\begin{equation}\label{eq:mpe_img_txt_impl}
\mathcal{X}^{\mathrm{img}} = g_0^{\mathrm{img}} + g_4^{\mathrm{img}}I,\qquad
\mathcal{X}^{\mathrm{txt}} = g_0^{\mathrm{txt}} + g_4^{\mathrm{txt}}I.
\end{equation}
因此它们被建模为 $\mathrm{SE}(3)$ 不变量语义通道，所有模态在 MPE 后直接拼接而不在嵌入阶段做额外融合。

\subsection{几何 Clifford 注意力}\label{sec:gca}
GCA 对每个阶分量分别施加独立线性映射，定义 $\mathcal{Q}_i=\mathcal{W}_Q(\mathcal{X}_i)$、$\mathcal{K}_j=\mathcal{W}_K(\mathcal{X}_j)$、$\mathcal{V}_j=\mathcal{W}_V(\mathcal{X}_j)$，并使用标量打分
\begin{equation}\label{eq:attn_score_impl}
a_{ij}=\frac{1}{\sqrt{d_{\mathcal{M}}}}\left\langle \mathcal{Q}_i\widetilde{\mathcal{K}}_j \right\rangle_0.
\end{equation}
标量权重不会改变值分量的表示类型，因此 GCA 的严格几何核心在于标量打分与标量加权聚合本身，而不在于任意按阶参数化。

\subsection{几何层归一化}\label{sec:gln}
GLN 按阶分离地使用
\begin{equation}\label{eq:gln_norm_impl}
\rho_r(\mathcal{X})=\sqrt{\left|\left\langle \langle\mathcal{X}\rangle_r\widetilde{\langle\mathcal{X}\rangle_r} \right\rangle_0\right|+\epsilon}
\end{equation}
进行归一化，并只对标量与伪标量引入固定偏置：
\begin{equation}\label{eq:gln_output_impl}
\operatorname{GLN}(\mathcal{X})=\sum_{r=0}^{4}\gamma_r\frac{\langle\mathcal{X}\rangle_r}{\rho_r(\mathcal{X})}+\beta_0+\beta_4I.
\end{equation}
由于协变分量的尺度在平移下并不严格保持，GLN 对完整 $\mathrm{SE}(3)$ 只能视为近似兼容。

\subsection{软等变正则化}\label{sec:equiv_reg}
训练阶段采样随机 motor $\mathcal{M}_g=T(\mathbf{t}_g)R(\boldsymbol{\omega}_g)$，并约束 $\Phi(\mathcal{M}_g\!\cdot\!\mathcal{X})$ 与 $\mathcal{M}_g\!\cdot\!\Phi(\mathcal{X})$ 在潜在多重向量空间中尽可能一致。距离函数采用按阶加权的二次型
\begin{equation}\label{eq:mv_distance_impl}
D(\mathcal{A},\mathcal{B}) = \sum_{r=0}^{4}\eta_r\left|\left\langle \langle\mathcal{A}-\mathcal{B}\rangle_r\widetilde{\langle\mathcal{A}-\mathcal{B}\rangle_r} \right\rangle_0\right|.
\end{equation}
当前代码默认采用各阶单位权重，并在计算该正则项时临时将模型切换到 $\mathrm{eval}$ 模式，以避免 dropout 噪声污染几何一致性度量。最终训练目标为 $\mathcal{L}_{\mathrm{total}}=\mathcal{L}_{\mathrm{task}}+\lambda\mathcal{L}_{\mathrm{equiv}}$，其中 $\lambda$ 采用线性预热。
